\documentclass[11pt]{article}
\usepackage{jotlocal}
\usepackage[numbers,sort&compress]{natbib}

\jottitle{The Assertion Ladder: Graduated Assertion for Embedding a Technical Domain in Repositories Shared with Coding Agents}

\jotauthors{Leigh Griffin$^{1}$ \and Paul Power$^{1}$ \and Ray Carroll$^{1}$ \and Colm Dunphy$^{2}$\\[0.3em]
\small$^{1}$Red Hat\qquad $^{2}$South East Technological University (SETU)}

\jotabstract{Coding agents now contribute to repositories alongside people. They write quickly,
they start every session with no memory of the last one, and their changes arrive faster than a
person can review them closely. Two forces follow: \emph{sprawl}, in which locally plausible
changes spread one domain rule across many places, and \emph{cognitive erosion}, in which
neither the reviewer nor the agent can hold enough of the domain in mind to notice. We argue that
a repository protects its domain to the extent that each domain rule is asserted at the right
pedagogical level, and we name the resulting pattern the \emph{Assertion Ladder}. Its six rungs
run from a rule that is merely implied, through one that is shown, told, checked and explained,
to one that cannot be written wrongly at all. Nine smaller patterns realise the rungs. Read
through lean thinking, the ladder is error-proofing and standard work for a contributor who
cannot remember. We present three receivers of the pattern, one small Kanban service implemented
identically in Java, Python and TypeScript, and show how each binds every rung in its own idiom.
Seeding fourteen canonical mistakes into each receiver, the single verify command catches all 42;
before the requirement and behavioural rungs were added, it caught 30 of 36 structural mistakes
and missed the same two in every language. A pilot with eighteen coding agents suggests that the
declared and enforced rungs reduce convention drift. The repository is public, and each receiver
can be read on its own.}

\jotkeywords{coding agents; software architecture; architecture testing; behaviour-driven
development; EARS; lean software development; design patterns; human--AI collaboration}

\begin{document}
\jotmaketitle

\section{Introduction}\label{sec:intro}

A new kind of contributor has joined software teams. Coding agents, language models given a
shell, a file system and a test runner, now read a repository, plan a change, edit many files
and run the tests before a person sees the result
\citep{chen2021codex,jimenez2024swebench,yang2024sweagent}. They are fast and broadly competent.
They are also, in one respect, the most difficult colleague a team has ever had: every session
starts with no memory of the last one. Whatever the agent learned about the domain yesterday,
about why the work-in-progress limit is checked in the service rather than the route, or why the
domain objects never import the web framework, is gone this morning.

The only thing a person and an agent reliably share is the repository. Whatever the repository
does not say, and does not enforce, the agent will reinvent, usually plausibly and often
differently each time. People have the opposite problem. They remember the domain, but they
cannot read every line an agent writes. Studies of programmers working with code-generating
tools describe a shift from writing to reviewing, and reviewing under that volume slides
towards skimming \citep{vaithilingam2022expectation,barke2023grounded}. Security studies show
what can pass unnoticed \citep{pearce2022asleep,perry2023insecure}.

We give the two resulting forces names. \emph{Sprawl} is what happens when a large volume of
locally plausible changes spreads one domain rule across many places. A check is copied into a
route because it was nearby. A notification is called inline because that was quicker than
publishing an event. A framework annotation lands on a domain object because it was convenient.
Every one of those changes would pass review on its own. Taken together they leave the domain
in pieces. \emph{Cognitive erosion} is what happens when neither party can hold enough of the
domain in mind to notice: the reviewer checks syntax because the intent is not visible in the
diff, and the agent re-derives conventions from scratch every session. Lehman observed long ago
that a system's structure decays under change unless work is done to preserve it
\citep{lehman1980laws}. Agents have not repealed that law. They have made it run faster.

\paragraph{Where we meet the problem} This paper comes out of a collaboration between a
university and an industrial engineering site in the same city, and we meet the problem from
both sides. At South East Technological University (SETU), students on the online Higher
Diploma in Science in Computer Science\footnote{\url{https://www.setu.ie/courses/higher-diploma-in-computer-science-2-years-online}} now learn to
program alongside coding agents, and their lecturers need codebases that teach a domain to both.
SETU's Master of Business Studies in Lean Enterprise
Excellence\footnote{\url{https://www.setu.ie/courses/master-of-business-in-lean-enterprise-excellence}} teaches the lean practice this
paper draws on. SETU delivers course material through Tutors, an open-source learning
platform whose production TypeScript codebase coding agents were already changing
\citep{tutorsmonorepo}. At Red Hat, the same agents contribute to engineering work, and hiring
is moving towards graduates who can work safely in repositories that agents also change. The
question is the same for a student, a new hire and an agent: how does a repository make its
domain rules impossible to miss and hard to break, without burying everyone in prose?

\paragraph{The need for constraints} The common response is to write more guidance: longer
prompts, longer contributor guides, longer instruction files for agents
\citep{agentsmd}. Guidance helps, but we have watched it fail in two ways. It is read at the
start of a session and then competes with everything else in the agent's context, and it grows
until nobody, person or agent, reads it at all. Our own first step beyond guidance was a
repository in which conventions are carried three ways at once: exemplified by one small
vertical slice, declared in a manifest of a few hundred words, and enforced by architecture
rules written as tests. The receivers in this paper began from that arrangement. This paper
generalises it. It asks not only \emph{whether} a rule is written down but \emph{how strongly}
it is asserted, and it extends the question from structure to behaviour and requirements.

\paragraph{The Assertion Ladder} Our central claim is a teaching one. A coding agent is a
capable learner with no memory, and the repository is its teacher. Teachers do not assert
everything the same way. They show a worked example, they tell a rule, they set an exercise that
marks the answer, they give feedback that explains the mistake, and now and then they arrange
the task so that the wrong answer is simply not available. The \emph{Assertion Ladder}
(Figure~\ref{fig:ladder}) gives each of these moves a rung, from \rung{L0 Implied} to
\rung{L5 Prevented}, and asks whoever designs a repository to place every domain rule on a rung
on purpose: high enough that the learner cannot miss it, and no higher than the rule's risk
justifies, because every assertion costs something to write, to maintain and to read.

\paragraph{Lean as the explanation} None of this is new engineering. It is old engineering
applied to a new worker. Lean thinking starts from \emph{value}, what the customer and, here,
the domain need, and treats everything that does not serve it as waste
\citep{womack1996lean,poppendieck2006implementing}. For work that repeats, its central tool is
\emph{standard work}: the current best method, followed by everyone who does the job and
improved by the team that owns it \citep{liker2004toyota}. It reads a process as a \emph{value
stream} and asks where along it a defect is found and what the waiting and rework cost
\citep{rother1999learning}. And it answers each problem with \emph{root cause analysis}, asking
why until the answer is something the process can change \citep{ohno1988tps}. Read this way,
the rungs are standard work for a contributor who cannot remember it: shown in a worked slice,
told in a short manifest, and then built into the work as checks that stop the build, explain
the stop, or make the mistake impossible to write. Section~\ref{sec:discussion} makes the case
that lean explains why the ladder works, rather than decorating it.

\paragraph{Critical gaps} Software engineering already has most of the parts. Architecture
conformance checking, specification by example, structured requirements and instruction files
for agents all exist, and Section~\ref{sec:related} reviews them. What it lacks is a standard
way to combine them. We identify five gaps in current practice that call for standardisation.
There is no shared vocabulary for how strongly a rule is asserted (\ref{gap:vocabulary}).
Instruction files for agents are unvalidated prose with no link to enforcement
(\ref{gap:prose}). Conformance tools say nothing about which rules to check rather than
document, or about a failure naming its fix (\ref{gap:selection}). Agent-facing repositories
have no executable closure between requirements and the scenarios that test them
(\ref{gap:closure}). And there is no benchmark of whether a repository's own harness catches
the mistakes that matter, as opposed to whether an agent can complete a task
(\ref{gap:benchmark}). None of these gaps is large on its own. Their combined effect is that
every team sharing a repository with agents invents its own answer, and most answer with more
prose.

\paragraph{Two known uses} We show the pattern in two settings. The first is greenfield:
one application, \emph{taskboard}, a small Kanban service with a work-in-progress limit, built
with the ladder from the start. We walk through its TypeScript implementation and show that
Java and Python implementations with identical behaviour bind every rung equivalently
\citep{taskboardrepo}. We call these implementations \emph{receivers} of the pattern rather
than its source, and work backwards from what each gains to the single capability they share.
The second is a retrofit: \emph{Tutors}, an open-source TypeScript learning platform in
production, to which the rungs were added after the fact while coding agents were already
contributing, and from which the release checks were then extracted into a reusable,
product-independent harness \citep{tutorsmonorepo,tutorsreleaseharness}. One shows the ladder
built where nothing existed. The other shows it fitted to a codebase with years of history.

\paragraph{Contributions} The paper makes five contributions, each aimed at one or more of
the gaps.
\begin{enumerate}
  \item The Assertion Ladder, a pattern in the classic form \citep{alexander1977pattern,gamma1994design}
    with six rungs and nine rung patterns, proposed as a shared vocabulary for how strongly a
    repository asserts each domain rule (\ref{gap:vocabulary}, \ref{gap:prose},
    \ref{gap:selection}; Section~\ref{sec:ladder}).
  \item A unification of three kinds of intent, structure, behaviour and requirements, as three
    loads on the same ladder, with a closure check that ties requirements to scenarios and the
    observation that the pattern of an EARS requirement \citep{mavin2009ears} predicts the layer
    that should implement it (\ref{gap:closure}; Sections~\ref{sec:ladder}
    and~\ref{sec:receivers}).
  \item A greenfield receiver in TypeScript, with equivalent Java and Python receivers, showing
    which rungs each language can reach and what it uses where it cannot
    (Section~\ref{sec:receivers}).
  \item A retrofit of a production codebase, Tutors, and the extraction of its release checks
    into a reusable harness, showing how the ladder is adopted where it was not designed in
    (Section~\ref{sec:tutors}).
  \item Violation seeding as a benchmark of a repository rather than of an agent, run before and
    after the behavioural and requirement rungs were added, and a pilot with eighteen coding
    agents (\ref{gap:benchmark}; Section~\ref{sec:evaluation}).
\end{enumerate}

\begin{figure*}[t]
  \centering
  \resizebox{\linewidth}{!}{% The Assertion Ladder: six rungs, each with its teaching move and its lean counterpart.
\begin{tikzpicture}[
    font=\small,
    rung/.style={draw=#1, line width=0.9pt, fill=#1!8, rounded corners=2pt,
                 minimum width=5.9cm, minimum height=0.78cm, anchor=west, align=left},
    note/.style={anchor=west, align=left, text width=4.6cm, font=\footnotesize},
  ]
  \def\gap{0.95}
  \foreach \i/\name/\teach/\lean/\col in {
      0/{L0 Implied}/{hidden curriculum}/{tacit knowledge, no standard}/muted,
      1/{L1 Shown}/{worked example}/{standard work, shown}/structure,
      2/{L2 Told}/{direct instruction}/{standard work, written}/requirement,
      3/{L3 Checked}/{assessment}/{built-in test at the point of work}/behaviour,
      4/{L4 Explained}/{formative feedback}/{andon: the stop says where and why}/reqc,
      5/{L5 Prevented}/{constraint}/{mistake-proofing (poka-yoke)}/structure} {
    \node[rung=\col] (r\i) at (0, \i*\gap) {\textbf{\name}\quad\textcolor{muted}{\footnotesize\teach}};
    \node[note] at (6.2, \i*\gap) {\textcolor{muted}{lean:} \lean};
  }
  % the rails of the ladder
  \draw[line width=1.6pt, jotblue] (-0.25,-0.5) -- (-0.25, 5*\gap+0.5);
  % arrows
  \draw[-{Stealth[length=3mm]}, line width=1pt, jotblue] (-0.8,-0.3) -- node[rotate=90, above, font=\footnotesize] {earlier, more certain feedback} (-0.8, 5*\gap+0.3);
  \draw[-{Stealth[length=3mm]}, line width=1pt, muted] (11.1, 5*\gap+0.3) -- node[rotate=-90, above, font=\footnotesize] {cheaper to build and maintain} (11.1,-0.3);
\end{tikzpicture}
}
  \caption{The Assertion Ladder. Each rung is a teaching move and a lean mechanism. Rules climb
  only as far as their risk warrants; as a rule reaches a checked rung, it is removed from the
  prose below it.}
  \label{fig:ladder}
\end{figure*}

\section{Background}\label{sec:background}

The Assertion Ladder draws on five bodies of work: the emergence of coding agents, architectural
structure and its testing, behaviour and requirement specification, the psychology of learning,
and lean production. This section introduces only what the later sections use.

\subsection{Coding agents and the memoryless contributor}

Code-generating models moved from completing a line \citep{chen2021codex} to resolving whole
issues in real repositories \citep{jimenez2024swebench}. What made the second step possible was
less the model than its harness: an agent that can list files, open them, edit them and run the
tests behaves very differently from one that emits a single completion
\citep{yang2024sweagent}. Tool vendors now converge on a convention for telling such an agent
how a repository works: a Markdown file at the root, loaded at the start of every session
\citep{agentsmd}. That file is the agent's whole inheritance from previous sessions.

Three properties of this contributor matter for what follows. It is \emph{capable}: given a
clear pattern it will reproduce it across layers and languages. It is \emph{memoryless}: nothing
learned in one session survives to the next unless the repository records it. And it is
\emph{fast}: its output can exceed what a person can review line by line, which moves the
person's role from author to reviewer \citep{barke2023grounded} and makes the quality of what
passes review depend on what the reviewer can see.

\subsection{Structure and its enforcement}

Layering is one of the oldest structuring ideas in software \citep{dijkstra1968the}; its
rationale is information hiding, so that a decision likely to change is confined to one module
\citep{parnas1972criteria}. The layered architecture of enterprise applications, with a
presentation layer, a service layer, a data-access layer and a domain model, is catalogued as a
pattern \citep{buschmann1996posa,fowler2002peaa}, and domain-driven design places the domain
model at the centre of attention \citep{evans2003ddd}. Hexagonal and Clean Architecture invert
the dependencies so that the domain depends on nothing and outer adapters depend on ports it
owns \citep{cockburn2005hexagonal,martin2017clean}.

Structure described only in prose decays. Architecture fitness functions make an
architectural characteristic executable, so that a build can fail when it is violated
\citep{ford2017evolutionary}; ArchUnit does this for Java dependencies by analysing bytecode in
an ordinary unit test \citep{archunit}. Such rules are the structural core of our third rung.

\subsection{Behaviour and requirements}

Behaviour-driven development writes tests as examples of behaviour in the vocabulary of the
domain, conventionally in the form \emph{given} a context, \emph{when} an event occurs,
\emph{then} an outcome is expected \citep{north2006bdd}; specification by example treats such
examples as the living specification of the system \citep{adzic2011sbe}.

The Easy Approach to Requirements Syntax (EARS) constrains natural-language requirements to a
small set of templates, each distinguished by a keyword \citep{mavin2009ears}:
\emph{ubiquitous} requirements (``The system shall\ldots''), \emph{state-driven}
(``While\ldots''), \emph{event-driven} (``When\ldots''), \emph{optional feature}
(``Where\ldots''), \emph{unwanted behaviour} (``If\ldots, then\ldots'') and \emph{complex}
combinations of these. EARS was designed to remove ambiguity for human readers. Section~\ref{sec:ladder}
argues that its keywords do more work in a repository shared with agents: they classify a rule
in a way that predicts where its code belongs.

\subsection{Learning: worked examples, scaffolding and cognitive load}

Cognitive load theory distinguishes the load a task imposes because of its intrinsic difficulty
from load imposed by how the task is presented, and holds that working memory is small
\citep{sweller1988cognitive}. One of its best-established findings is the worked-example
effect: novices learn a procedure faster by studying a solved example than by solving the
problem themselves \citep{sweller1985worked}. Scaffolding describes support that lets a learner
do what they could not do alone, and that is withdrawn as competence grows \citep{wood1976scaffolding}.

A coding agent is not a novice in the usual sense; it is highly competent in general and
entirely ignorant of this repository in particular. The analogy we use is limited to that
second property. The repository must teach the particular, every session, within a context
window that is large but not unlimited.

\subsection{Lean production and lean software}

Lean production, derived from the Toyota Production System, treats anything that does not add
value for the customer as waste, and builds quality into the process rather than inspecting it
in at the end \citep{ohno1988tps,womack1996lean}. Four of its mechanisms recur in this paper.
\emph{Jidoka} stops the line at the first defect so that it cannot propagate.
\emph{Poka-yoke} designs a fixture or step so that a mistake cannot be made, or is detected at
the moment it is made \citep{shingo1986zqc}. \emph{Standard work} records the current best method
so that it can be followed and improved. \emph{Andon} gives every worker a way to stop and call
for help. Lean software development maps these onto programming practices such as automated
tests, continuous integration and small batches \citep{poppendieck2003lean}, and the
build--measure--learn loop brings the scientific method to product decisions
\citep{ries2011lean}.

\subsection{Patterns}

A pattern names a recurring problem in a context, the forces that make it hard, and a solution
that balances them \citep{alexander1977pattern}. The form was brought to software design by the
Gang of Four \citep{gamma1994design} and to architecture by the POSA series
\citep{buschmann1996posa}. We use it for the Assertion Ladder because the ladder is not an
algorithm or a tool but a way of resolving forces that every team sharing a repository with
agents will meet.

\section{The Assertion Ladder}\label{sec:ladder}

This section states the pattern, describes each rung, names the smaller patterns that realise
the rungs, and shows that structure, behaviour and requirements are three loads carried by the
same ladder rather than three separate techniques.

\subsection{The pattern}

\begin{description}[leftmargin=2.2cm,style=nextline]
  \item[Name] Assertion Ladder (also: Graduated Assertion).
  \item[Intent] Keep a technical domain intact when coding agents contribute, by asserting each
    domain rule at the weakest level that still covers its risk.
  \item[Context] A repository shared by people and coding agents that start every session with
    no memory. Changes arrive faster than anyone can review them closely.
  \item[Problem] How do a domain's rules survive contributors who cannot remember them and
    reviewers who cannot read everything?
  \item[Forces] See Figure~\ref{fig:forces}.
    \begin{itemize}
      \item A rule nobody wrote down is lost at the next session.
      \item Written guidance grows until nobody reads it, and competes with everything else in
        an agent's context.
      \item Every executable check costs effort to write, time in the build, and maintenance
        when the domain changes.
      \item The strongest assertions, those made by the type system, depend on the language.
      \item Too much constraint makes the repository hard to change. Ease of use is a force, not
        an afterthought.
      \item A check that an agent can edit is a check that an agent can weaken.
    \end{itemize}
  \item[Solution] For each domain rule, choose a rung. Climb until the rule's risk is covered,
    then stop. When a rule reaches a checked rung, remove it from the prose so that the
    declaration stays short. Run every check behind one command, and put the checks themselves
    behind a human boundary.
  \item[Consequences] Feedback arrives early and names the fix; the manifest stays short;
    intent can be reviewed separately from code. In exchange there are more tests to maintain,
    the top rung differs between languages, and a rule that climbs further than its risk
    warrants makes change expensive.
\end{description}

\begin{figure*}[t]
  \centering
  \resizebox{\linewidth}{!}{% Two forces acting on the domain, and the ladder standing between them and the domain.
\begin{tikzpicture}[
    font=\small,
    box/.style={draw=#1, line width=0.9pt, fill=#1!7, rounded corners=3pt, align=center, inner sep=6pt},
    arr/.style={-{Stealth[length=2.6mm]}, line width=0.9pt, #1},
  ]
  \node[box=jotblue, minimum width=3.2cm, minimum height=1.6cm] (domain)
    {\textbf{The domain}\\[2pt]\footnotesize rules in one place,\\\footnotesize in its own words};
  \node[box=missed, left=2.6cm of domain, text width=3.6cm] (sprawl)
    {\textbf{Sprawl}\\[2pt]\footnotesize locally plausible changes\\\footnotesize spread one rule\\\footnotesize across many places};
  \node[box=behaviour, right=2.6cm of domain, text width=3.6cm] (erosion)
    {\textbf{Cognitive erosion}\\[2pt]\footnotesize the reviewer cannot read it all;\\\footnotesize the agent cannot remember};
  \node[box=structure, below=1.1cm of domain, text width=8.2cm] (ladder)
    {\textbf{Assertion Ladder}: move each rule out of memory and prose\\into the repository's own checks, as high as its risk warrants};
  \draw[arr=missed] (sprawl) -- node[above, font=\footnotesize]{dilutes} (domain);
  \draw[arr=behaviour] (erosion) -- node[above, font=\footnotesize]{hides} (domain);
  \draw[arr=structure] (ladder) -- node[right, font=\footnotesize]{protects} (domain);
  \draw[arr=structure, dashed] (ladder.west) -| node[pos=0.75, left, font=\footnotesize]{fails fast} (sprawl.south);
  \draw[arr=structure, dashed] (ladder.east) -| node[pos=0.75, right, font=\footnotesize]{shows intent} (erosion.south);
\end{tikzpicture}
}
  \caption{Two forces act on a domain once agents contribute. Sprawl spreads a rule across
  places; cognitive erosion hides the spread from both reviewer and agent. The ladder answers
  both by moving rules out of memory and prose into the repository's own checks.}
  \label{fig:forces}
\end{figure*}

\subsection{The rungs}

\paragraph{\rung{L0 Implied}} The rule exists only in someone's head or in habit. The agent
experiences nothing, and guesses. The guess is usually plausible, which is what makes L0
dangerous: nothing about the result looks wrong. Lean would call this tacit knowledge with no
standard: nothing to follow, nothing to improve, and the cost of relearning paid every session.
A teacher would call it the hidden curriculum.

\paragraph{\rung{L1 Shown}} One complete example embodies the rule: a vertical slice from the
domain to the endpoint, small enough to read whole. Agents, like novices, copy structure before
they understand it \citep{sweller1985worked}, so the slice is the cheapest effective assertion
there is. To a lean eye the slice is standard work shown rather than written: the team's current
best way to add a feature, in a form a contributor can copy. Its weakness is that the learner may
copy the wrong detail, or generalise from an accident of the example.

\paragraph{\rung{L2 Told}} A short declaration states the rule, in a file the agent's tooling
loads at the start of every session \citep{agentsmd}. Prose can say what code cannot: what needs
a person's decision, which of two acceptable choices the team prefers, why a rule exists. It is
the written half of standard work. Its weakness is that it can be forgotten, misread or crowded
out, and that it grows.

\paragraph{\rung{L3 Checked}} An executable rule fails the build when the rule is broken. This
is the rung of architecture tests, contract tests and behavioural scenarios. The agent
experiences a red suite and must work out why. This is the built-in test that
Spear and Bowen found attached to Toyota's standard work \citep{spear1999decoding}: the defect is
found where it is made, by the contributor who made it, before the change moves on to a reviewer,
rather than by a separate inspection further down the stream.

\paragraph{\rung{L4 Explained}} The check's failure names the rule, the offender and the fix. A
bare assertion failure tells the agent \emph{that} it is wrong; an explained failure tells it
\emph{what to do instead}, which is the difference between summative and formative feedback. A
bare failure is a stop. An explained failure is a stop that, like an andon signal, shows where
the problem is and why, and names the countermeasure. It costs a sentence per rule.

\paragraph{\rung{L5 Prevented}} The wrong form does not compile. A sealed hierarchy with an
exhaustive switch, or a map keyed by a closed union type, makes an unmapped error impossible to
write. This is mistake-proofing (poka-yoke) of the preventive kind: the fixture admits the part
only one way, so the mistake is never made, where a check at L3 catches it once it has been. It is
also the rung that depends most on the language.

\subsection{Nine rung patterns}

Each rung is realised by one or more smaller patterns, listed in Table~\ref{tab:rung-patterns}.
Two are not tied to a rung: \pattern{Single Gate}, which runs every check behind one command so
that ``finished'' has one meaning, and \pattern{Human Boundary}, which lists the few changes an
agent may not make alone, including changes to the checks themselves.

\begin{table*}[t]
  \centering\small
  \caption{The rung patterns. Each answers one recurring problem.}
  \label{tab:rung-patterns}
  \begin{tabularx}{\linewidth}{l|l|X|X}
    \hline
    \textbf{Rung} & \textbf{Pattern} & \textbf{Problem} & \textbf{Solution} \\
    \hline\hline
    L1 & \pattern{Worked Slice} & How does a newcomer learn the shape of a feature? & One complete vertical slice, small enough to read whole. \\
    \hline
    L2 & \pattern{Short Manifest} & How is what cannot be tested passed on? & One file of a few hundred words, loaded every session. \\
    \hline
    L3 & \pattern{Executable Rule} & How is a structural rule kept when nobody remembers it? & A test that fails when the rule is broken. \\
    \hline
    L3 & \pattern{Closure Check} & How do two lists stay in step? & Assert that every X has a Y. \\
    \hline
    L3 & \pattern{Honest Double} & How do fast tests stay truthful? & One contract, run against the real implementation and its stand-in. \\
    \hline
    L4 & \pattern{Teaching Failure} & How does a red build lead to the right fix? & The failure names the rule, the offender and the remedy. \\
    \hline
    L5 & \pattern{Unrepresentable Mistake} & How is a mistake stopped before any test runs? & Types that make the wrong form fail to compile. \\
    \hline
    all & \pattern{Single Gate} & When is work finished? & One command runs every rung; green means done. \\
    \hline
    all & \pattern{Human Boundary} & Who may change the ladder itself? & A short list of changes that need a person first. \\
    \hline
  \end{tabularx}
\end{table*}

\pattern{Closure Check} recurs more than any other pattern, so it is worth a paragraph of its
own. An
\emph{every X has a Y} assertion keeps two lists in step without either list knowing about the
other: every domain error has an HTTP status; every repository implementation passes the same
contract; every requirement has a scenario and every scenario cites a requirement. It is the
cheapest way we know to make a growing set safe to extend, and extending sets is most of what an
agent does all day.

\subsection{Three kinds of intent, one ladder}

Architecture tests, behaviour-driven development and structured requirements are usually
presented as separate practices with separate literatures. In a repository shared with agents
they are the same move applied to three kinds of intent
(Table~\ref{tab:intents}). \emph{Structure} says where code may live and what it may touch.
\emph{Behaviour} says what the system does, in the domain's own words. \emph{Requirements} say
why, once, in a form that can be traced. Each is declared, exemplified and enforced, and each
has its own closure check.

\begin{table*}[t]
  \centering\small
  \caption{The same three carriers hold three kinds of intent. Text colour follows the intent
  throughout the paper.}
  \label{tab:intents}
  \begin{tabularx}{\linewidth}{>{\raggedright}p{2.6cm}|X|X|X}
    \hline
    \textbf{Intent} & \textbf{Declared (L2)} & \textbf{Exemplified (L1)} & \textbf{Enforced (L3--L5)} \\
    \hline\hline
    \textcolor{structure}{\textbf{Structure}}: where code may live &
      a placement table in the manifest &
      the move-task slice, domain to route &
      architecture rules; exhaustive error mapping \\
    \hline
    \textcolor{behaviour}{\textbf{Behaviour}}: what the system does &
      one convention line: scenarios cite ids and read given / when / then &
      the work-in-progress scenario &
      the scenarios, and the repository contract \\
    \hline
    \textcolor{requirement}{\textbf{Requirements}}: why &
      \repo{REQUIREMENTS.md}, one EARS sentence per rule &
      ten taskboard rules with ids &
      the requirement closure check \\
    \hline
  \end{tabularx}
\end{table*}

\subsection{EARS as a placement oracle}

The most useful property of EARS in this setting was not one it was designed for. Its keywords
classify a rule, and the class predicts where the rule's code belongs (Table~\ref{tab:ears}).
An \emph{unwanted behaviour} about one entity (``If a requested move is not allowed from the
task's status, then\ldots'') is a guard on that entity that raises a domain error. A
\emph{state-driven} rule that spans many entities (``While the work-in-progress limit is
reached\ldots'') needs a collaborator and belongs in a service. An \emph{event-driven} rule whose
response is to tell someone (``When a task moves to done, the board shall announce\ldots'')
is a domain event published by the service with the effect in a handler. An \emph{optional
feature} (``Where a limit is configured\ldots'') is a field on the one settings object. A
\emph{ubiquitous} rule (``The board shall answer every refused request with a problem
document'') is a cross-cutting invariant: one place, guarded by a check.

\begin{table*}[t]
  \centering\small
  \caption{From EARS pattern to layer, with the taskboard rule that exercises each row.}
  \label{tab:ears}
  \begin{tabularx}{\linewidth}{l|X|X}
    \hline
    \textbf{EARS pattern} & \textbf{Implement as} & \textbf{Taskboard rule} \\
    \hline\hline
    Unwanted, one entity & a guard on the entity raising a domain error & R-FLOW-2: an invalid move is refused \\
    \hline
    State-driven, many entities & a service method with a collaborator & R-WIP-1: the work-in-progress limit \\
    \hline
    Event-driven, tells someone & a domain event; the effect in a handler & R-DONE-1: completion is announced \\
    \hline
    Optional feature & a field on the settings object & the configurable limit \\
    \hline
    Ubiquitous & one place, guarded by a closure or architecture check & R-ERR-1: every refusal is a problem document \\
    \hline
  \end{tabularx}
\end{table*}

The mapping matters because the commonest placement mistake an agent makes is to put a correct
rule in the wrong place, and a behavioural test cannot see that mistake (Section~\ref{sec:evaluation}).
A requirement written in EARS carries its placement with it. We present the mapping as a
hypothesis supported by the ten taskboard rules and by one pilot run, not as a law.

\subsection{Compression and fading}

The solution says to climb only as far as a rule's risk warrants, and to remove a rule from the
prose once a check covers it. Both clauses are about ease of use, which is the force teams
forget first. A manifest that repeats what the tests enforce is waste in the lean sense: it costs
reading time every session and adds nothing the gate does not already guarantee. Standard work
is meant to change as the team learns, and here it changes by subtraction as well as addition:
once a rule is built into a check, the written standard no longer needs to carry it. Instructional
design calls the same move \emph{fading}: scaffolding is withdrawn as the learner no longer needs
it \citep{wood1976scaffolding}. In the receivers the manifest says so of itself: ``anything a
test can check lives in the tests, not here.'' Section~\ref{sec:evaluation} reports what the
added rungs cost in words.

\section{Three receivers}\label{sec:receivers}

A pattern is only as credible as its known uses. This section presents three: one application,
\emph{taskboard}, implemented three times with identical behaviour. We call them
\emph{receivers} because they do not define the Assertion Ladder; they take it on, each in the
idiom of its language, and each gains something from it. The paper works backwards from those
gains to the capability they share.

\subsection{The application}

Taskboard is a Kanban API with one entity. A task moves from \emph{todo} to \emph{in progress}
to \emph{done}, and may move back from \emph{in progress} to \emph{todo}. No more than a
configured number of tasks may be in progress at once. When a task is completed the team is
told. Every refused request is answered with an RFC~9457 problem document carrying a stable
error code. The application is deliberately small: 283 lines of application code in Java, 223
in Python and 242 in TypeScript, down from 637, 729 and 957 in a first iteration that used
dependency-injection containers, object-relational mappers and a second entity
\citep{griffin2026convention}. Small enough to read whole is a property the ladder depends on:
the \pattern{Worked Slice} only works if the slice can be read in one sitting.

Each receiver is organised in the same six layers: \repo{api}, \repo{services},
\repo{repositories}, \repo{domain}, and the side layers \repo{events} and \repo{config}, with one
composition root. Imports point only downwards, the domain imports nothing but the language's
standard library, and each outside capability (HTTP, SQL, the environment, logging) belongs to
exactly one layer.

\subsection{A steel thread through all three kinds of intent}

Figure~\ref{fig:thread} follows one rule, the work-in-progress limit, from its requirement to
the gate. The requirement is written once in EARS:

\begin{quote}\small
\textbf{R-WIP-1} (state-driven). While \texttt{WIP\_LIMIT} tasks are in progress, when another
task is moved to \texttt{in\_progress}, the board shall refuse with
\texttt{wip\_limit\_exceeded} and leave the task where it was.
\end{quote}

\noindent Its pattern predicts a service method, because the rule needs to count other tasks.
The scenario that covers it cites its id and reads as given / when / then
(Listing~\ref{lst:scenario}). The domain error it raises carries a code and no HTTP; one handler
maps it to a status, and a closure check refuses a new error without one. The limit arrives
through a constructor from the one settings object. Every step is checked by the same command.

\begin{figure}[t]
  \centering
  % One rule, the work-in-progress limit, through every kind of intent and every layer.
\begin{tikzpicture}[
    font=\footnotesize,
    step/.style={draw=#1, line width=0.9pt, fill=#1!7, rounded corners=2pt, text width=3.05cm,
                 minimum height=1.35cm, align=center, inner sep=4pt},
    arr/.style={-{Stealth[length=2.3mm]}, line width=0.8pt, muted},
  ]
  \node[step=requirement] (req) {\textbf{Requirement}\\R-WIP-1, state-driven\\\texttt{REQUIREMENTS.md}};
  \node[step=behaviour, right=0.45cm of req] (sc) {\textbf{Scenario}\\given a full board\ldots\\cites R-WIP-1};
  \node[step=structure, right=0.45cm of sc] (svc) {\textbf{Service}\\guard in \texttt{move\_task}\\(needs other tasks)};
  \node[step=structure, right=0.45cm of svc] (err) {\textbf{Domain error}\\\texttt{WipLimitExceeded}\\code, no HTTP};
  \node[step=structure, below=0.9cm of err] (map) {\textbf{One translation}\\error $\to$ 409\\exhaustive map};
  \node[step=structure, left=0.45cm of map] (cfg) {\textbf{Setting}\\\texttt{WIP\_LIMIT}\\read once, injected};
  \node[step=muted, left=0.45cm of cfg, text width=6.55cm] (gate) {\textbf{Single gate}: compiler, architecture rules,\\requirement check, scenarios, contract, API tests};
  \draw[arr] (req) -- (sc);
  \draw[arr] (sc) -- (svc);
  \draw[arr] (svc) -- (err);
  \draw[arr] (err) -- (map);
  \draw[arr] (svc.south) |- ($(cfg.north)+(0,0.3)$) -- (cfg.north);
  \draw[arr] (map) -- (cfg);
  \draw[arr] (cfg) -- (gate);
\end{tikzpicture}

  \caption{The work-in-progress limit as a steel thread. Colours mark the kind of intent each
  step carries; every step is checked by the single gate.}
  \label{fig:thread}
\end{figure}

\begin{lstlisting}[language=Python,caption={A scenario in the Python receiver. The docstring cites the requirement and reads given / when / then; the requirement check fails if either is missing.},label=lst:scenario]
def test_the_wip_limit_is_enforced(service: TaskService) -> None:
    """R-WIP-1: given a full board, when another task is started,
    then it is refused and stays in todo."""
    first, second = service.create_task("First"), service.create_task("Second")
    service.move_task(first.id, Status.IN_PROGRESS)
    with pytest.raises(WipLimitExceeded):
        service.move_task(second.id, Status.IN_PROGRESS)
    assert service.get_task(second.id).status is Status.TODO
\end{lstlisting}

\subsection{Binding paths}

Table~\ref{tab:binding} shows how each receiver binds every rung pattern. Reading a row shows
one pattern bound three ways; reading a column shows one receiver's complete path.

\begin{table}[t]
  \centering\small
  \caption{Binding paths of the three receivers.}
  \label{tab:binding}
  \begin{tabularx}{\linewidth}{@{}>{\raggedright\arraybackslash}p{3.1cm} X X X@{}}
    \toprule
    Rung pattern & Java (Spring Boot) & Python (FastAPI) & TypeScript (Fastify) \\
    \midrule
    \pattern{Worked Slice} & \repo{moveTask}, service to controller & \repo{move\_task}, service to route & \repo{moveTask}, service to route \\
    \pattern{Short Manifest} & \repo{AGENTS.md}, 480 words & \repo{AGENTS.md}, 427 words & \repo{AGENTS.md}, 468 words \\
    \pattern{Executable Rule} & ArchUnit over bytecode, 8 rules & AST walk, 7 rules incl.\ strict type check & source text, 5 rules \\
    \pattern{Closure Check} & errors by the compiler; requirements by \repo{RequirementsTest} & errors and requirements by tests & errors by the compiler; requirements by a test \\
    \pattern{Honest Double} & contract run on JDBC and in-memory & contract run on SQLite and in-memory & contract run on SQLite and in-memory \\
    \pattern{Teaching Failure} & \repo{because(\ldots)} on every rule & a message on every assertion & a message on every assertion \\
    \pattern{Unrepresentable Mistake} & sealed \repo{DomainException}, exhaustive \repo{switch} & not available; bound at L3 & \repo{Record<ErrorCode, number>} \\
    \pattern{Single Gate} & \repo{mvn verify} & \repo{pytest} & \repo{npm test} \\
    \pattern{Human Boundary} & \multicolumn{3}{>{\raggedright\arraybackslash}p{\dimexpr\linewidth-3.8cm\relax}}{the same three items in each manifest: dependencies; the rules and manifest; schema and endpoint shape} \\
    \bottomrule
  \end{tabularx}
\end{table}

\paragraph{Java} The Java receiver reaches the top of the ladder most often. Its domain errors
form a sealed hierarchy, and the handler that maps them to HTTP uses a \texttt{switch} with no
default branch, so a new error does not compile until it has a status. Its architecture rules
are ArchUnit rules over bytecode, the most precise of the three enforcement styles, and each
carries a \texttt{because} clause that ArchUnit prints on failure (Listing~\ref{lst:archunit}).
Its requirement check uses ArchUnit to import the test classes themselves and asserts that every
scenario's \texttt{@DisplayName} opens with a requirement id.

\begin{lstlisting}[language=Java,caption={An executable rule with a teaching failure in the Java receiver. It closes a gap that no behavioural test can see (Section~\ref{sec:evaluation}).},label=lst:archunit]
@ArchTest
static final ArchRule onlyEventHandlersHaveSideEffects = noClasses()
    .that().resideOutsideOfPackage("..events..")
    .should().dependOnClassesThat().resideInAnyPackage(
        "org.slf4j..", "java.util.logging..", "org.apache.commons.logging..")
    .because("side effects (logs, mail, audit) listen for a domain "
           + "event; publish one instead");
\end{lstlisting}

\paragraph{Python} Python cannot make an unmapped error unrepresentable, so the Python receiver
binds that rule one rung lower: a \pattern{Closure Check} asserts that every subclass of
\texttt{DomainError} has an entry in the status map. It compensates elsewhere. Its architecture
rules walk the syntax tree, which lets them express rules the other receivers cannot, such as
``a route contains no \texttt{if}, loop, comprehension, \texttt{raise} or \texttt{try}''. And the
strict type checker runs inside the same gate, so type errors are reported as failures of the
single command rather than of a separate tool.

\paragraph{TypeScript} The TypeScript receiver reaches L5 through the type system: the status
map is a \texttt{Record} keyed by a closed union of error codes, so a missing entry is a
compile error (Listing~\ref{lst:record}). Its architecture rules are the cheapest of the three,
regular expressions over source text, and the least precise; they are adequate for a codebase
this small and would need a parser for a larger one.

\begin{lstlisting}[caption={Unrepresentable Mistake in the TypeScript receiver.},label=lst:record]
// Record<ErrorCode, number> is exhaustive: a new error code
// will not compile until it has a status.
const STATUS_BY_CODE: Record<ErrorCode, number> = {
  task_not_found: 404,
  invalid_transition: 409,
  wip_limit_exceeded: 409,
};
\end{lstlisting}

\subsection{Layered, not Clean, and what that costs}

All three receivers are layered rather than hexagonal or Clean. The persistence interface lives
beside its implementation in the repositories layer rather than as a port owned by the core, so
dependency inversion does not stop a service from reaching SQL. Capability confinement does:
a rule that only the repositories layer may use the database driver, only the api layer the web
framework, only the config layer the environment, and only the events layer logging. We read
this as a finding about the ladder rather than about architecture styles. A handful of
\pattern{Executable Rule}s recovers most of the domain protection Clean Architecture provides,
without the ports, adapters and interactors an agent would otherwise have to learn and
reproduce in every feature. The trade is ceremony against rules; the ladder lets a team make it
explicitly.

\subsection{What each receiver gains}

Working backwards from the receivers, each gains the same four things. The domain stays in one
place, which the seeding results in Section~\ref{sec:evaluation} measure directly. A new feature
has a recipe: state the rule in EARS, let its pattern choose the layer, copy the slice, run the
gate. Failures say what to do. And the manifest a new contributor reads is under 500 words,
because everything checkable lives in the checks. The receivers differ only in which rung
carries which rule, and those differences are themselves instructive: the same rule set,
bound by different languages, lands on different rungs.

\section{Use cases}\label{sec:usecases}

The receivers show the pattern bound in code. This section states, independently of any
language, the situations in which a domain is at risk once agents contribute, so that the
pattern can be applied to domains other than task boards. Each use case names a situation, the
threat, the invariant at stake and the teaching move that answers it. Table~\ref{tab:usecases}
lists twelve in six families and Figure~\ref{fig:coverage} places each on the ladder, both in the
receivers as first published and as they stand after the work described here.

\begin{table}[t]
  \centering\small
  \caption{Twelve use cases in six families. The last two columns give the rung at which the
  receivers asserted the invariant at commit \texttt{8c4a74a} and now.}
  \label{tab:usecases}
  \begin{tabularx}{\linewidth}{@{}l X X c c@{}}
    \toprule
    & Situation & Invariant at stake & Then & Now \\
    \midrule
    \multicolumn{5}{@{}l}{\textit{Placement: sprawl spreads one rule across many places}} \\
    UC1 & A memoryless agent adds a feature & every part lands in the layer that owns it & L3 & L4 \\
    UC2 & A side effect is added to a use case & use cases publish facts; reactions subscribe & L2 & L4 \\
    UC3 & A business rule drifts to the edge & routes translate; they never decide & L0 & L4 \\
    \multicolumn{5}{@{}l}{\textit{Contamination: convenience pulls infrastructure into the core}} \\
    UC4 & A framework type enters the domain & the domain imports only the standard library & L3 & L4 \\
    UC5 & An outside capability leaks into a service & each capability has one owning layer & L3 & L4 \\
    \multicolumn{5}{@{}l}{\textit{Contract integrity: the promises the domain makes}} \\
    UC6 & A new failure mode is introduced & every error has one public translation & L5 & L5 \\
    UC7 & A test double lies & every double passes the real contract & L3 & L3 \\
    \multicolumn{5}{@{}l}{\textit{Harness erosion: the protections themselves wear down}} \\
    UC8 & The agent edits a rule to go green & rules change only with a person's decision & L2 & L2 \\
    UC9 & A dependency is added & new dependencies need a person & L2 & L2 \\
    \multicolumn{5}{@{}l}{\textit{Cognitive load: what people and agents can hold in mind}} \\
    UC10 & A reviewer faces a large AI diff & intent is reviewable apart from code & L1 & L3 \\
    UC11 & An agent starts cold & the declaration stays small & L3 & L3 \\
    \multicolumn{5}{@{}l}{\textit{Transfer: taking the approach to a new domain}} \\
    UC12 & A new domain is embedded & each invariant is placed on the ladder deliberately & L1 & L1 \\
    \bottomrule
  \end{tabularx}
\end{table}

\begin{figure}[t]
  \centering
  % Use cases on the ladder: hollow circle at 8c4a74a, filled circle now, arrow where they moved.
\begin{tikzpicture}[font=\footnotesize, x=0.95cm, y=0.52cm]
  \foreach \l in {0,...,5} {
    \draw[muted!35] (\l,0.4) -- (\l,-12.4);
    \node[above, font=\footnotesize\bfseries, text=jotblue] at (\l,0.4) {L\l};
  }
  \foreach \y/\name/\then/\now in {
      1/{UC1 feature by a memoryless agent}/3/4,
      2/{UC2 side effect in a use case}/2/4,
      3/{UC3 business rule at the edge}/0/4,
      4/{UC4 framework in the domain}/3/4,
      5/{UC5 capability leaks}/3/4,
      6/{UC6 new failure mode}/5/5,
      7/{UC7 lying test double}/3/3,
      8/{UC8 agent edits a rule}/2/2,
      9/{UC9 dependency added}/2/2,
      10/{UC10 reviewer overload}/1/3,
      11/{UC11 cold start}/3/3,
      12/{UC12 new domain}/1/1} {
    \node[anchor=east] at (-0.5,-\y) {\name};
    \ifnum\then=\now\relax
      \fill[structure] (\now,-\y) circle (4.2pt);
    \else
      \draw[-{Stealth[length=2mm]}, structure, line width=0.8pt] (\then+0.18,-\y) -- (\now-0.2,-\y);
      \draw[structure, line width=0.9pt, fill=white] (\then,-\y) circle (4.2pt);
      \fill[structure] (\now,-\y) circle (4.2pt);
    \fi
  }
\end{tikzpicture}

  \caption{Where each use case sits on the ladder in the receivers, at commit \texttt{8c4a74a}
  (hollow) and now (filled). UC6 is at L5 in Java and TypeScript and at L3 in Python.}
  \label{fig:coverage}
\end{figure}

\subsection{Placement}

Placement failures are the signature of sprawl. In \textbf{UC1} an agent with no prior session
implements a feature end to end; the risk is not that the feature fails but that each layer gets
a plausible and inconsistent piece of it. The answer is a worked slice to copy, followed by
architecture rules to pass. \textbf{UC2} and \textbf{UC3} are subtler because the mistake
preserves behaviour. A notification called inline in a service, instead of published as an
event, produces the same log line; a work-in-progress check copied into a route produces the
same refusal. No behavioural test can tell the difference, and at commit \texttt{8c4a74a}
nothing in any receiver did. They were closed by two \pattern{Executable Rule}s, that only the
events layer may log or notify and that routes make no decisions, each with a
\pattern{Teaching Failure} naming where the code should go.

\subsection{Contamination}

In \textbf{UC4} a framework annotation looks convenient on a domain object; in \textbf{UC5} a
service needs a setting or a quick query. Both are answered by capability confinement
(Section~\ref{sec:receivers}), and both were already checked before this work; they gained
explained failures.

\subsection{Contract integrity}

\textbf{UC6}, a new failure mode, is the one use case the receivers answer at L5 in two of three
languages: the new error does not compile until it has a status. \textbf{UC7} is the test double
that lies, answered by an \pattern{Honest Double}. It has a history: in the first seeding run the
contract created tasks titled First, Second and Third, whose alphabetical order equals their
creation order, so a double that sorted by title could not be caught. Changing the titles to
Write, Review and Publish fixed it. The episode is a small instance of kaizen, and a reminder that
a check is only as good as its fixture.

\subsection{Harness erosion}

\textbf{UC8} and \textbf{UC9} concern the ladder itself. An agent whose build is red has an easy
way to make it green: relax the rule. No agent in our pilot did so, but only the manifest forbade
it (L2). Raising UC8 to L3 is a matter of repository settings rather than tests, for example by
requiring a named person's review for any change to the rule files. Dependencies (UC9) are
similar, and can be checked by pinning the dependency list in a test.

\subsection{Cognitive load}

\textbf{UC10} is the reviewer's problem. When a change adds a feature, the parts a person most
needs to judge are the new rule and the examples that illustrate it, not the plumbing. Keeping
requirements in their own file, with scenarios that cite them and a check that keeps the two in
step, puts the intent of a change in a small, separate diff. We claim this lowers the reviewer's
load; we have not yet measured it. \textbf{UC11} is the agent's problem: every session starts
cold, and guidance that grows to compensate ends up unread. Fading keeps the manifest under 500
words even after two new kinds of intent were added.

\subsection{Transfer}

\textbf{UC12} asks whether the pattern survives a move to a domain that is not a task board. Two
illustrations, not implemented, suggest how. In a double-entry ledger, the invariant that debits
equal credits can be stated in EARS as an unwanted behaviour, shown in a posting slice, checked by
a property test, and prevented by a transaction type that can only be built from balanced legs.
In medication dosing, the invariant that a dose may not exceed a weight-based maximum without a
clinician's override can be checked by boundary scenarios and prevented, in part, by types that
distinguish milligrams from milligrams per kilogram. In both, the question the ladder asks is the
same: how high must this rule climb, given what it costs to get it wrong?

\section{Evaluation}\label{sec:evaluation}

We ask three questions. \textbf{RQ1}: which mistakes does each rung catch, and which does it
miss? \textbf{RQ2}: what do the added rungs cost in size and reading? \textbf{RQ3}: do the
declared and enforced rungs change what coding agents do? Every number below can be regenerated
from the repository's \repo{evaluation/} directory \citep{taskboardrepo}.

\subsection{RQ1: violation seeding}

A seeding script plants one canonical mistake at a time in a receiver, runs its single verify
command, records the earliest mechanism that objects, and restores the file. It refuses to start
unless the receiver is green, and checks that it is green again at the end. The mechanisms, in
order of how early they give feedback, are the compiler, an architecture rule, the requirement
check, a behavioural test, and nothing. Table~\ref{tab:violations} lists the fourteen mistakes:
twelve structural, and two that break the trace between requirements and scenarios.

\begin{table*}[t]
  \centering\small
  \caption{The fourteen seeded mistakes, and the mechanism that caught each in Java / Python /
  TypeScript now. C compiler, A architecture rule, R requirement check, B behavioural test.}
  \label{tab:violations}
  \begin{tabularx}{\linewidth}{l|X|c}
    \hline
    & \textbf{Mistake} & \textbf{J / P / T} \\
    \hline\hline
    V1 & The API bypasses the service and uses a repository & A / A / A \\
    \hline
    V2 & A service signals failure with an HTTP concept & A / A / B \\
    \hline
    V3 & A service issues SQL & A / A / A \\
    \hline
    V4 & A component constructs its own collaborator & A / A / A \\
    \hline
    V5 & A service reads the environment & A / A / A \\
    \hline
    V6 & A new domain error has no HTTP status & C / A / C \\
    \hline
    V7 & A route catches a domain error and answers ad hoc & B / A / A \\
    \hline
    V8 & A side effect is inlined in the service & A / A / A \\
    \hline
    V9 & The test double drifts from the real repository & B / B / B \\
    \hline
    V10 & A business rule is placed in the route & A / A / A \\
    \hline
    V11 & A setting is hard-coded in the service & B / B / B \\
    \hline
    V12 & A framework type enters the domain & A / A / A \\
    \hline
    V13 & A requirement has no scenario & R / R / R \\
    \hline
    V14 & A scenario cites no requirement & R / R / R \\
    \hline
  \end{tabularx}
\end{table*}

Figure~\ref{fig:seeding} compares the receivers at commit \texttt{8c4a74a}, before the work
described here, with the receivers now. Before, the single gate caught 30 of 36 structural
mistakes. The six misses were the same two mistakes in every language: V8, a side effect inlined
in the service, and V10, a business rule placed in the route. That the gap was identical across
three languages and three enforcement styles is itself a result: it belonged to the capability,
not to any receiver. Both mistakes preserve behaviour, so no behavioural test could have caught
them; only a structural rule could. After two such rules were added, and the requirement check
with its two seeds, the gate catches all 42 mistakes. The new routes rule also moved V7 from a
behavioural test to an architecture rule in Python and TypeScript, which is earlier feedback for
the same mistake. Application code did not change between the two runs; only tests, manifests
and requirement files did.

\begin{figure*}[t]
  \centering
  % Seeding results, before (8c4a74a, 12 structural seeds) and now (14 seeds), per receiver.
\begin{tikzpicture}
  \begin{axis}[
      xbar stacked,
      width=0.86\linewidth,
      height=6.2cm,
      bar width=9pt,
      xmin=0, xmax=14,
      xtick={0,2,...,14},
      xlabel={seeded mistakes},
      ytick=data,
      yticklabels={{TypeScript, now},{TypeScript, then},{Python, now},{Python, then},{Java, now},{Java, then}},
      yticklabel style={font=\footnotesize},
      tick label style={font=\footnotesize},
      label style={font=\footnotesize},
      legend style={at={(0.5,-0.24)}, anchor=north, legend columns=5, font=\footnotesize, draw=none},
      enlarge y limits=0.12,
      axis x line*=bottom,
      axis y line*=left,
      xmajorgrids, grid style={muted!20},
    ]
    % order of rows (bottom to top): Java then, Java now, Python then, Python now, TS then, TS now
    \addplot[fill=compiler, draw=none] coordinates {(1,5) (1,4) (0,3) (0,2) (1,1) (1,0)};
    \addplot[fill=caught, draw=none]   coordinates {(6,5) (8,4) (7,3) (10,2) (5,1) (8,0)};
    \addplot[fill=reqc, draw=none]     coordinates {(0,5) (2,4) (0,3) (2,2) (0,1) (2,0)};
    \addplot[fill=behav, draw=none]    coordinates {(3,5) (3,4) (3,3) (2,2) (4,1) (3,0)};
    \addplot[fill=missed, draw=none]   coordinates {(2,5) (0,4) (2,3) (0,2) (2,1) (0,0)};
    \legend{compiler, architecture rule, requirement check, behavioural test, not caught}
  \end{axis}
\end{tikzpicture}

  \caption{Seeded mistakes by the earliest mechanism that caught them, per receiver, at commit
  \texttt{8c4a74a} (12 structural seeds) and now (14 seeds).}
  \label{fig:seeding}
\end{figure*}

\subsection{RQ2: cost}

Table~\ref{tab:size} reports size. Application code is unchanged. Test code grew by between 31
and 89 lines per receiver, and each manifest by between 80 and 83 words, remaining under the 500
words the plan allows. The added rungs therefore cost about one screen of test code and one
paragraph of prose per receiver, most of it the requirement check and the ids on existing tests.

\begin{table*}[t]
  \centering\small
  \caption{Size of the receivers at commit \texttt{8c4a74a} and now (\repo{evaluation/metrics*.json}).}
  \label{tab:size}
  \begin{tabular}{l|r|r|r|r|r|r}
    \hline
    \multicolumn{1}{l|}{} & \multicolumn{2}{c|}{\textbf{Application lines}} & \multicolumn{2}{c|}{\textbf{Test lines}} & \multicolumn{2}{c}{\textbf{Manifest words}} \\
    \cline{2-3}\cline{4-5}\cline{6-7}
    \textbf{Receiver} & \textbf{then} & \textbf{now} & \textbf{then} & \textbf{now} & \textbf{then} & \textbf{now} \\
    \hline\hline
    Java       & 283 & 283 & 287 & 376 & 397 & 480 \\
    \hline
    Python     & 223 & 223 & 193 & 254 & 347 & 427 \\
    \hline
    TypeScript & 242 & 242 & 193 & 224 & 388 & 468 \\
    \hline
  \end{tabular}
\end{table*}

\subsection{RQ3: a pilot with coding agents}

Eighteen coding agents, each starting with no memory, implemented the same feature request, task
deletion, in copies of the receivers at commit \texttt{8c4a74a}. Each copy offered one, two or
three of the carriers: the exemplar only (L1), the exemplar and the manifest (L1--L2), or the
repository as published (L1--L3). The design crossed three languages, three conditions and two
models, with one run per cell. A hidden acceptance test of four cases, never shown to an agent,
judged the result, and a scripted audit fixed in advance of inspection recorded where each agent
placed each part of the change.

\begin{table*}[t]
  \centering\small
  \caption{Pilot results by condition (six runs each; \repo{evaluation/agent\_experiment/}).}
  \label{tab:pilot}
  \begin{tabular}{l|c|c|c|c}
    \hline
    \textbf{Condition} & \textbf{Acceptance passed} & \textbf{Departures from convention} & \textbf{Mean tokens} & \textbf{Mean seconds} \\
    \hline\hline
    Exemplar (L1)            & 6 / 6 & 2 & 61\,650 & 121 \\
    \hline
    + Manifest (L1--L2)      & 6 / 6 & 0 & 61\,990 & 132 \\
    \hline
    + Enforced (L1--L3)      & 6 / 6 & 0 & 65\,696 & 151 \\
    \hline
  \end{tabular}
\end{table*}

All eighteen runs delivered a passing suite and passed every acceptance case
(Table~\ref{tab:pilot}). Two departed from the conventions, both by the smaller model in the
exemplar-only condition. One, in Java, logged the deletion inside the service instead of
publishing an event: the V8 mistake, which no rule caught at the time. The deletion request's
last clause, ``whenever a task is deleted the team must be told'', is an event-driven requirement
in EARS terms, and the mapping in Table~\ref{tab:ears} places it in an event handler. The other,
in Python, did not add the new operation to the in-memory double, which only the type checker in
the full rule set notices. With one run per cell these results are suggestive, not significant.
The pilot has not yet been repeated against the receivers as they stand now, with the deletion
request written in EARS; that is the next experiment.

\section{Discussion}\label{sec:discussion}

\subsection{Lean explains the ladder}

Lean thinking faced a version of the problem this paper faces: work done quickly by many hands,
more of it than anyone could inspect, where a defect that travelled downstream cost far more to
fix than one caught where it was made \citep{ohno1988tps,womack1996lean}. Its answer was to build
the specification into the work itself. We read the Assertion Ladder as that answer applied to
a repository and a memoryless contributor, and this section makes the case through four Lean
ideas: value, standard work, the value stream and root cause analysis. Table~\ref{tab:Lean} maps each to the
receivers and to the evidence, and the section closes with the places where the analogy stops.

\begin{table*}[t]
  \centering\small
  \caption{Lean ideas, the rungs and patterns that carry them, and where they appear in the
  receivers. The first four rows are the argument; the last three are supporting mechanisms.}
  \label{tab:Lean}
  \begin{tabularx}{\linewidth}{>{\raggedright\arraybackslash}p{2.5cm}|>{\raggedright\arraybackslash}p{2.9cm}|X|X}
    \hline
    \textbf{Lean idea} & \textbf{Rung or pattern} & \textbf{In the receivers} & \textbf{Evidence} \\
    \hline\hline
    Value, and waste removed & fading; \pattern{Short Manifest} & no DI container, no ORM, one entity; no prose for what a test checks & 637 / 729 / 957 $\to$ 283 / 223 / 242 lines; manifests 480 / 427 / 468 words \\
    \hline
    Standard work & L1 \pattern{Worked Slice}, L2 \pattern{Short Manifest}, with L3 as its built-in test & the slice and the feature recipe, followed by every contributor and changed only by people & every pilot run followed the slice's layering \\
    \hline
    Value stream: defects found earlier & L3--L5 & V8 and V10 now stopped at the gate, not left to review; V7 moved from a behavioural test to an architecture rule in Python and TypeScript & V8 and V10 passed the gate in all three receivers at \texttt{8c4a74a}; caught now \\
    \hline
    Root cause analysis and kaizen & \pattern{Executable Rule}, \pattern{Teaching Failure}; \pattern{Honest Double} & the misses traced to one cause and closed by two rules; the V9 fixture corrected & 30 / 36 $\to$ 42 / 42; V9 caught in all three receivers \\
    \hline
    Jidoka (the line stops itself) & \pattern{Single Gate} & one verify command; ``change the code, not the test'' & no pilot run modified a rule (18 / 18) \\
    \hline
    Andon (signal and cord) & L4 \pattern{Teaching Failure}; \pattern{Human Boundary} & failures name the rule and the fix; ``ask a human first'' for dependencies, rules and schema & message on every rule; boundary declared (L2), not yet enforced \\
    \hline
    Poka-yoke (mistake-proofing) & L5 \pattern{Unrepresentable Mistake} & exhaustive error mapping that will not compile when a case is missing & V6 caught by the compiler in Java and TypeScript \\
    \hline
  \end{tabularx}
\end{table*}

\paragraph{Value} Lean asks first what the customer values, and the ladder adds the domain to the
answer: a change is valuable only if it works and leaves the domain's rules in one place. Sprawl is
then easy to name as waste. It is work that looks finished but will be redone, and the rework
falls on whoever next tries to change the rule. The same lens explains the force most easily
forgotten, ease of use. A paragraph of prose that repeats what a test enforces is read by every
contributor in every session and adds nothing the gate does not already guarantee; for a
contributor that relearns everything each time, that reading is the dominant cost. Fading removes
it. The receivers show both halves: application code fell from 637, 729 and 957 lines to 283, 223
and 242 when containers, mappers and a second entity were removed, and the manifests stayed under
500 words while two new kinds of intent were added.

\paragraph{Standard work} Standard work is where the mapping is tightest. The worked slice and the
feature recipe of Section~\ref{sec:receivers} (state the rule in EARS, let its pattern choose the
layer, copy the slice, run the gate) specify the content of a change, its sequence and its
expected outcome, which is how Spear and Bowen describe Toyota's work \citep{spear1999decoding}.
The ladder then does what they found Toyota does: it attaches a test to the standard, so that a
departure shows up at once. The lower rungs show and tell the standard; the upper rungs make it
self-checking. Every pilot run followed the slice's layering, and the two departures the pilot
recorded both came from the condition with the slice alone and no written standard. Standard work
is also owned. In Lean, the people who do the work improve the standard; in the receivers, an agent
follows the standard, and the \pattern{Human Boundary} requires a person's decision before the
rules or the manifest change. That division is deliberate, and we return to it below.

\paragraph{The value stream} Figure~\ref{fig:valuestream} draws the value stream of one change:
written by an agent, verified by the single gate, reviewed by a person, merged and released. Each
rung decides where along that stream a broken rule is found. At L0 only a reviewer who reads
closely can see it, and cognitive erosion is the loss of that reading. At L1 and L2 the
standard makes the mistake less likely but finds nothing. At L3 and L4 the gate finds it before
any person has spent time on the change, and at L5 it cannot be written. The seeding results show
the shift. At commit \texttt{8c4a74a}, V8 and V10 passed the gate in all three receivers and would
have reached review; now both are stopped at the gate by an architecture rule. V7 moved earlier
within the gate, from a behavioural test to an architecture rule, in Python and TypeScript, and V6
was already refused by the compiler in Java and TypeScript. We have no timing data for the stream, so we claim
only the position at which a defect is found, not the time saved. The Lean argument is that
position matters: a defect found at the gate goes back to the contributor who made it while the
change is still open, whereas one found in review has already cost a person's attention and, if
missed, becomes the rework that sprawl describes.

\begin{figure*}[t]
  \centering
  \resizebox{0.92\linewidth}{!}{% The value stream of one change, and where each rung finds a defect along it.
\begin{tikzpicture}[
    font=\footnotesize,
    step/.style={draw=#1, line width=0.9pt, fill=#1!7, rounded corners=2pt, text width=2.2cm,
                 minimum height=1.1cm, align=center, inner sep=4pt},
    arr/.style={-{Stealth[length=2.3mm]}, line width=0.8pt, #1},
    note/.style={text width=2.3cm, align=center, font=\scriptsize, anchor=north},
  ]
  % the stream
  \node[step=structure] (write) {\textbf{Write}\\an agent edits};
  \node[step=behaviour, right=0.55cm of write] (verify) {\textbf{Verify}\\the single gate};
  \node[step=missed, right=0.55cm of verify] (review) {\textbf{Review}\\a person reads};
  \node[step=muted, right=0.55cm of review] (merge) {\textbf{Merge}};
  \node[step=muted, right=0.55cm of merge] (prod) {\textbf{Release}};
  \draw[arr=muted] (write) -- (verify);
  \draw[arr=muted] (verify) -- (review);
  \draw[arr=muted] (review) -- (merge);
  \draw[arr=muted] (merge) -- (prod);
  % where each rung acts
  \node[note] at ($(write.south)+(0,-0.15)$) {L1, L2 guide the change\\L5 refuses to compile it\\[2pt]\textcolor{muted}{V6 in Java and TypeScript}};
  \node[note] at ($(verify.south)+(0,-0.15)$) {L3 finds the defect\\L4 names the fix\\[2pt]\textcolor{muted}{V8, V10 now}};
  \node[note] at ($(review.south)+(0,-0.15)$) {L0: only a person\\can see it\\[2pt]\textcolor{muted}{V8, V10 at \texttt{8c4a74a}}};
  \node[note] at ($(merge.south)+(0,-0.15)$) {nothing new\\is checked};
  \node[note] at ($(prod.south)+(0,-0.15)$) {found by a user,\\or never};
  % rework loops
  \draw[arr=behaviour] (verify.north) -- ++(0,0.35) -| node[pos=0.25, above]{rework at once; nobody waits} ($(write.north)+(0.45,0)$);
  \draw[arr=missed] (review.north) -- ++(0,0.95) -| node[pos=0.25, above]{rework after a person has waited} ($(write.north)+(-0.45,0)$);
  % direction of the ladder
  \draw[arr=jotblue, line width=1pt] ($(prod.south)+(0,-1.55)$) -- node[below]{each rung climbed moves detection upstream} ($(write.south)+(0,-1.55)$);
\end{tikzpicture}
}
  \caption{The value stream of one change. Each rung moves the point at which a broken rule is
  found further upstream. At commit \texttt{8c4a74a}, V8 and V10 passed the gate and could be seen
  only in review; they are now stopped at the gate. The figure shows position only; we have no
  timing data.}
  \label{fig:valuestream}
\end{figure*}

\paragraph{Root cause analysis} The first seeding run is a small worked example of Lean problem
solving. The problem was concrete: six of 36 seeded mistakes passed the gate. Asking why led to a
cause at the level of the ladder, not of any receiver. The misses were the same two mistakes, V8
and V10, in all three languages. Both passed because every behavioural test passed; every
behavioural test passed because neither mistake changes what the system does; and nothing else
looked, because the structural rules covered which layer may import which, not which layer may
log, notify or decide. The root cause was that behaviour cannot see placement and no rule asserted
placement for side effects or decisions. The countermeasure was two \pattern{Executable Rule}s,
each with a \pattern{Teaching Failure} naming where the code belongs, which is to say two rules
moved from L0 and L2 to L4. The check was a second run, which caught all 42 seeds, including the new requirement
seeds. The standard then
absorbed the countermeasure, in the plan--do--check--act manner of an A3 report
\citep{sobek2008a3}. The V9 fixture episode (Section~\ref{sec:usecases}) followed the same cycle on
a smaller scale, and is the kaizen of a check rather than of the code. What Lean adds to the usual
advice to add a test after a bug is the insistence on the cause: because the same gap appeared in
three languages, the countermeasure was a kind of rule every receiver needed, not a patch to one of
them.

Figure~\ref{fig:a3} writes the episode up as an A3, in the form the Tutors release harness
writes for a release (Section~\ref{sec:tutors}). We include it because it is the ladder's
hallmark on one page. The current condition is where a mistake was found along the stream, the
root cause is the rung it was sitting on, the countermeasure is the rung it climbed to, and the
check is the next run. The form carries no verdict of its own, which is the point. In Tutors the
harness opens an A3 on every failed candidate and its findings inform a confidence score that a
person reads before shipping. In the receivers the countermeasure went straight into the gate.
A team can use the same sheet for confidence or for a hard check, and the ladder is the scale
on which it says which.

\begin{figure*}[t]
  \centering
  \resizebox{0.96\linewidth}{!}{% The first seeding run written up as an A3 report, in the form the Tutors harness writes for a release.
\begin{tikzpicture}[
    font=\scriptsize,
    panel/.style={draw=jotblue!60, line width=0.6pt, fill=white, text width=7.1cm, align=left,
                  inner sep=5pt, anchor=north west},
    head/.style={font=\scriptsize\bfseries, text=jotblue},
    step/.style={draw=#1, line width=0.7pt, fill=#1!7, rounded corners=2pt, minimum height=0.55cm,
                 minimum width=1.5cm, align=center, inner sep=2pt},
    arr/.style={-{Stealth[length=2mm]}, line width=0.7pt, muted},
  ]
  % sheet
  \draw[jotblue, line width=1pt, fill=wash] (0, 0) rectangle (15.4, -10.6);
  \node[anchor=north west, font=\footnotesize\bfseries, text=jotblue] at (0.25, -0.2)
    {A3: six seeded mistakes passed the gate};
  \node[anchor=north east, font=\scriptsize, text=muted] at (15.15, -0.25)
    {receivers at \texttt{8c4a74a}, 21 September 2026; countermeasures measured 29 September};

  % left column
  \node[panel] (bg) at (0.25, -0.8) {
    \textbf{\textcolor{jotblue}{Background}}\\
    The three receivers enforce structure with architecture rules and a short manifest.
    Seeding twelve canonical mistakes into each and running the single verify command,
    the gate caught 30 of 36.};
  \node[panel] (cur) at (0.25, -2.55) {
    \textbf{\textcolor{jotblue}{Current condition}}\\[3pt]
    \begin{tikzpicture}[baseline]
      \node[step=structure] (w) {write};
      \node[step=behaviour, right=0.5cm of w] (g) {gate};
      \node[step=missed, right=0.5cm of g] (r) {review};
      \draw[arr] (w) -- (g);
      \draw[arr] (g) -- node[above, text=missed, font=\tiny]{V8, V10 pass} (r);
    \end{tikzpicture}\\[2pt]
    V8 (a side effect inlined in the service) and V10 (a business rule placed in the route)
    pass every test in Java, Python and TypeScript and reach a human reviewer. The other
    thirty are stopped at the gate.};
  \node[panel] (goal) at (0.25, -5.0) {
    \textbf{\textcolor{jotblue}{Goal}}\\
    Every seeded mistake stopped at the gate, by the earliest mechanism its language allows,
    with no change to application code.};
  \node[panel] (rca) at (0.25, -6.3) {
    \textbf{\textcolor{jotblue}{Root cause (five whys)}}\\
    1. Why did V8 and V10 pass? Every behavioural test passed.\\
    2. Why? Neither mistake changes what the system does.\\
    3. Why did no rule object? The rules say which layer may import which, not which layer
       may log, notify or decide.\\
    4. Why? No rule asserted placement for side effects or decisions.\\
    5. Why in all three languages? Behaviour cannot see placement, and the pattern had no
       structural rule for it. The gap was in the ladder, not in a receiver.};

  % right column
  \node[panel] (cm) at (7.95, -0.8) {
    \textbf{\textcolor{jotblue}{Countermeasures}}\\
    Two \textsc{Executable Rule}s in each receiver: only the events layer may log or notify;
    routes decide nothing. Each carries a \textsc{Teaching Failure} that names where the
    code belongs. V8 and V10 move from L0 and L2 to L4.\\[2pt]
    While there: a requirement \textsc{Closure Check}, with two new seeds (V13, V14) to
    prove it can fail.};
  \node[panel] (chk) at (7.95, -3.45) {
    \textbf{\textcolor{jotblue}{Check}}\\
    Second run: 42 of 42 caught. V7 now found by an architecture rule rather than a
    behavioural test in Python and TypeScript. Cost: 31 to 89 lines of test code and
    80 to 83 words of manifest per receiver. Application code unchanged.};
  \node[panel] (fu) at (7.95, -5.55) {
    \textbf{\textcolor{jotblue}{Follow-up}}\\
    Repeat the pilot against the current receivers with the request in EARS. Raise the
    harness-erosion use cases from L2 to L3. Adopt the release harness on a second product.};
  \node[panel, fill=jotblue!6] (how) at (7.95, -7.3) {
    \textbf{\textcolor{jotblue}{How this sheet is used}}\\
    In Tutors the harness writes an A3 for every release and its findings inform a confidence
    score. In the receivers the countermeasure became part of the gate. The form is the same;
    a team chooses whether it advises or stops the line.};
\end{tikzpicture}
}
  \caption{The first seeding run as an A3 report. One sheet carries the problem, where it was
  found, why, what rung the rule climbed to and what the next run showed. Tutors writes the same
  sheet for a release and reads it as confidence; the receivers made its countermeasure a gate.}
  \label{fig:a3}
\end{figure*}

\paragraph{Supporting mechanisms} Three further Lean mechanisms map onto the ladder more briefly.
The single gate is \emph{jidoka}: the build stops itself on an abnormality, so that no person has
to watch every line, and the manifest's instruction to change the code rather than the test is
the rule that a stop is answered by fixing its cause. No pilot run modified a rule. Jidoka's
older name, automation with a human touch, fits a team of people and agents well: the machine
detects and stops, the person judges. The \pattern{Teaching Failure} is the \emph{andon} signal,
which shows where the problem is and why; the \pattern{Human Boundary} is the andon cord, the
right and the duty to stop and ask a person before changing a dependency, a rule or a schema. The
\pattern{Unrepresentable Mistake} is \emph{poka-yoke} of the preventive kind.

\paragraph{Where the analogy stops} The mapping is close but not exact, and the places where it
breaks tell us something. First, an agent does not learn across shifts. A factory worker is taught the
standard, practises it and in time no longer needs the sheet; an agent is taught afresh every
session. The cost of showing and telling is therefore paid every time, which is the strongest
reason to move a rule from prose into a check, where following it costs no reading. For the same
reason, improvement cannot happen inside the agent. Whatever an agent notices in a session is lost
unless it is written into the repository, so reflection and kaizen run through people and the
files they own. Second, Lean's respect for people becomes the \pattern{Human Boundary}. In Lean the
people who do the work own the standard. Here much of the work is done by agents, so people keep
what needs judgement: what counts as value, how high a rule should climb, and when a check is
wrong. The risk to avoid is the one Lean warned against, reducing people to inspectors of someone
else's output; the ladder gives them the standard and its checks to own instead of every line to
read. Third, the fixture is made of the same material as the work. A worker cannot redesign the jig
in the middle of a shift, but an agent can edit a test, which is why changes to the rules sit
behind the \pattern{Human Boundary}, and why use case UC8 remains at L2 until repository settings
enforce it. Fourth, the ladder can itself become waste. A check is cheap to run but not free to
maintain, and a rule encoded in types far more elaborate than its risk warrants is
over-processing. Reinertsen's caution applies \citep{reinertsen2009principles}, and so do the
findings of Staats and colleagues on uncertain knowledge work \citep{staats2011lean}: standardise
what should repeat and leave the rest free. The ladder applies Lean to a team's conventions, not
to its design.

\subsection{Domain sanctity against ease of use}

A repository can be made arbitrarily hard to damage by making it arbitrarily hard to change.
The ladder resists this in two ways. Each rule climbs only as far as its risk warrants, so the
cost of a rung is paid only where it buys something. Lean would call the alternative
over-processing. And the ladder is compressed from below as
it grows from above: when the receivers gained a requirement file, a scenario convention and
four new rules, their manifests grew by one paragraph each and stayed under 500 words, so the
written standard stayed short enough to be read in every session. The
evidence that ease of use survived is modest but concrete: in the pilot, agents given the full
rule set finished in 151 seconds on average against 121 for the exemplar alone, and used about
6\% more tokens, for a condition in which none departed from the conventions.

\subsection{Where each kind of intent stops}

Each kind of intent sees some mistakes and is blind to others, which is why the three are
complementary rather than redundant. Behaviour cannot see placement: V8 and V10 pass every
scenario because what the system does is unchanged. Structure cannot see meaning: an
architecture rule will not notice that a task may now skip a step. Requirements cannot execute:
without the closure check that ties them to scenarios, they decay into prose, the failure the
manifest's own design avoids. The ladder's contribution is to put all three on the same scale,
so that a team can ask of any rule, in any of the three, how strongly it is asserted.

\subsection{Threats to validity}

\paragraph{Construct} Seeded mistakes are canonical, not observed; real agents may make
mistakes the fourteen do not represent. The TypeScript architecture rules match source text, not
syntax, and could be evaded by code written to avoid them.

\paragraph{Internal} The authors built the receivers, the seeds and the rules together with an
AI coding agent, in the collaborative mode the paper studies. Rules were added after the misses
they close were known; the second seeding run therefore shows that the gap can be closed, not
that the rules generalise to unseen mistakes.

\paragraph{External} One small domain, three mainstream stacks and two models from one vendor
limit generality. The pilot ran one agent per cell and has not yet been repeated against the
current receivers.

\paragraph{Conclusion} Seeding results are deterministic and reproducible from the repository.
Pilot results are suggestive only.

\section{Related work}\label{sec:related}

\paragraph{Architecture conformance} Checking code against an intended architecture has a long
history; architecture fitness functions \citep{ford2017evolutionary} and tools such as ArchUnit
\citep{archunit} bring it into the ordinary test suite. The Assertion Ladder places such checks
at one rung among six and adds what they lack on their own: a view of when a rule should be
checked rather than shown or told, and a requirement that the check teach.

\paragraph{Specification by example and structured requirements} BDD \citep{north2006bdd} and
specification by example \citep{adzic2011sbe} make behaviour executable; EARS makes requirements
unambiguous \citep{mavin2009ears}. We combine them with a closure check that keeps requirements
and scenarios in step, and we observe a use for EARS it was not designed for: predicting
placement.

\paragraph{Architecture styles} Layered, hexagonal and Clean architectures
\citep{buschmann1996posa,fowler2002peaa,cockburn2005hexagonal,martin2017clean,evans2003ddd}
protect the domain by structure. The receivers show that, for a domain this size, capability
confinement rules give a layered codebase much of the same protection with less ceremony.

\paragraph{Coding agents} Benchmarks and agent harnesses \citep{chen2021codex,jimenez2024swebench,yang2024sweagent}
measure whether an agent can complete a task. Studies of programmers using code generators
\citep{vaithilingam2022expectation,barke2023grounded} and of the security of generated code
\citep{pearce2022asleep,perry2023insecure} show what passes review. Instruction files for agents
\citep{agentsmd} are the L2 rung in practice. Our concern is different: not whether one change
is correct, but whether the domain survives many correct-looking changes.

\paragraph{Learning and lean} The pedagogical framing borrows from cognitive load theory and the
worked-example effect \citep{sweller1988cognitive,sweller1985worked} and from scaffolding
\citep{wood1976scaffolding}; the lean framing from the Toyota Production System and its software
descendants \citep{ohno1988tps,shingo1986zqc,womack1996lean,poppendieck2003lean,ries2011lean}.

\paragraph{Our earlier work} The repository standard of exemplified, declared and enforced
conventions \citep{griffin2026convention} is the structural special case of the ladder. This
paper extends it to behaviour and requirements, grades assertion by strength, and reports the
second seeding run.

\todo{Related work needs a systematic search before submission; the entries above are a starting set.}

\section{Conclusion}\label{sec:conclusion}

Coding agents make a domain's rules easier to break and harder to notice breaking. The
Assertion Ladder answers by asking, for each rule, how strongly the repository asserts it: shown,
told, checked, explained or prevented. Structure, behaviour and requirements turn out to be three
loads on the same ladder, and the pattern of an EARS requirement suggests where its code belongs.
A greenfield receiver in TypeScript, with equivalent receivers in Java and Python, binds every
rung in its language's idiom; seeding fourteen canonical mistakes into each, the single verify
command catches all 42, where it caught 30 of 36 structural mistakes before the behavioural and
requirement rungs and two new structural rules were added. The cost was about a screen of tests
and a paragraph of prose per receiver. The Tutors retrofit shows the same rungs established in a
production codebase that agents were already changing, reaching the checked and explained rungs
wherever a check could be written, adding ratchets so that legacy violations could only shrink,
and extracting its release checks into a harness that the change under review cannot weaken.

We therefore propose the ladder as a standard for repositories shared with coding agents, in
answer to the five gaps of Section~\ref{sec:related}: a shared vocabulary of rungs; instruction
files that state only what is not yet checked; checks chosen by risk and written to teach;
requirements closed against scenarios; and violation seeding as a benchmark of the repository
itself rather than of the agent that works in it.

Lean thinking explains why this works. It counts as value only what the customer, and here the
domain, needs, and treats the rest as waste. It makes the best current method standard work that
everyone follows and the team improves, it asks where along the value stream a defect is found,
and it answers each miss by asking why until the process can change. The ladder is standard work,
with its tests built in, for a contributor who cannot remember it; each rung it climbs moves
detection further upstream; and the two seeding runs are one turn of lean's improvement cycle,
from a miss to its root cause to a countermeasure that was measured again.

Four directions follow. The pilot should be repeated against the current receivers, with more
runs per cell and the feature request written in EARS, to test whether stated requirements
improve placement. The harness-erosion use cases should be raised from L2 to L3 by protecting
the rule files themselves. The extracted release harness should be adopted by a second product,
to test how much of it is generic. And the reviewer's load, which we claim the ladder reduces by
separating intent from code, should be measured.

\paragraph{Availability} The receivers, the seeding script, the pilot's workspaces, audit and
diffs, and this paper's sources are public \citep{taskboardrepo}, as are Tutors and its release
harness \citep{tutorsmonorepo,tutorsreleaseharness}. Each receiver can be read on its own,
starting from its \repo{AGENTS.md}.


\bibliographystyle{plainnat}
\bibliography{references}

\clearpage
\appendix
\section{Reading a receiver}\label{app:repository}

Each receiver is self-contained. A reader interested in one language can start from its
directory and read in this order.

\begin{enumerate}
  \item \repo{AGENTS.md}: the manifest. Verify command, placement table, conventions,
    requirement convention, feature recipe, and what needs a person first.
  \item \repo{REQUIREMENTS.md}: the ten rules in EARS, with ids and the pattern-to-layer table.
  \item The worked slice: the domain entity, the service's move-task method, the repository
    interface and its implementation, and the route or controller.
  \item The scenarios: unit tests over the service with an in-memory repository, and API tests
    through the real composition root. Each cites the requirements it covers.
  \item The rules: the architecture tests and the requirement check.
\end{enumerate}

\begin{table}[!ht]
  \centering\small
  \caption{Where to find each part in each receiver.}
  \begin{tabularx}{\linewidth}{@{}l X X X@{}}
    \toprule
    Part & \repo{java/} & \repo{python/} & \repo{typescript/} \\
    \midrule
    Verify & \repo{mvn verify} & \repo{pytest} & \repo{npm test} \\
    Architecture rules & \repo{ArchitectureTest.java} & \repo{test\_architecture.py}$^{*}$ & \repo{architecture.test.ts}$^{*}$ \\
    Requirement check & \repo{RequirementsTest.java} & \repo{test\_requirements.py}$^{*}$ & \repo{requirements.test.ts}$^{*}$ \\
    Scenarios & \repo{service/}, \repo{api/}, \repo{events/} tests & \repo{tests/unit}, \repo{tests/api} & \repo{tests/unit}, \repo{tests/api} \\
    Honest double & \repo{TaskRepository\-Contract.java} & \repo{tests/integration/} & \repo{tests/integration/} \\
    \bottomrule
  \end{tabularx}\\[0.3em]
  {\footnotesize $^{*}$ under \repo{tests/architecture/}.}
\end{table}

The evaluation scripts are in \repo{evaluation/}: \repo{measure.py} for size,
\repo{seed\_violations.py} for seeding, and \repo{agent\_experiment/} for the pilot. The first
paper's results are kept alongside the current ones as \repo{*-8c4a74a.json}.


\end{document}
